\documentclass[12pt]{extarticle}
\usepackage[utf8]{inputenc}
\usepackage{hyperref}
\usepackage{enumitem}
\usepackage{parskip}
\usepackage{setspace}
\usepackage{titlesec}
\usepackage{amsmath}
\usepackage[margin=0.8in]{geometry}
\usepackage{fancyhdr}

\title{\Huge Groups, Series, and Sequences}
\author{\Large Keith Smith}
\date{}

\pagestyle{fancy}
\fancyhf{}
\fancyfoot[C]{\thepage}
\renewcommand{\headrulewidth}{0pt}

\begin{document}

\maketitle

\vspace{-2em}

\section*{Structure of the Visual Book}

\subsection*{TYPES of ORGANIZATIONS}

\begin{flushright}
\textit{Single pictures amassed are either a \emph{group,} \emph{series,} or a \emph{sequence.}}
\end{flushright}

A \textit{group} is a list. It is held in union by a common denominator, such as subject matter, composition, or from a similar source. A group is a \textit{compilation} without structure or constructed movement. Reference is not made from picture to picture; consequently, there is no set order of viewing. A group is topical and cannot be modified by storyline or directed movement.

Once a reference is intended within a group of pictures by the bookmaker, the unit is transformed with the characteristic of directed movement (referral) from one picture to another. The pictures can no longer be considered a group. They are now either a series or a sequence, depending upon how the bookmaker structures the movement. The book does not have to be solely a series or a sequence. A number of pictures, for instance, a chapter, might be ordered as a series. Then, several pictures might be presented as a group. The following chapter may be a sequence, or revert to a series.

\begin{flushright}
\textit{A group is a \emph{collection} without structure or constructed movement.}
\end{flushright}

A \textit{series} \textit{links} a number of pictures in a straight line; it is a \textit{contiguous} structure. A series is a linear, arithmetical progression. Each picture, an extension of the previous, modifies the next: a succession, a metamorphosis, or a narration. (It is important to note that definitions are tricky. In literature, a ``narrative story'' may be constructed as a series, but might be a sequence.)

A series is linking movement.

\textbf{A sequence} is constructed by cause and effect; it is contingent in structure. Several pictures react to, or act upon, each other, but not necessarily with the adjacent picture as in a series. The structure is contrapuntal. A sequence is a geometric progression, a montage. (It is important to note that the terminology in filmmaking is contradictory: a filmic ``sequence'' is a series by my definition.)

A sequence is conditional movement.

\textbf{CONFUSION of TERMS:} The concept of a sequence is complicated. It is compounded by the misuse of the term as a substitution for the verb \textit{order}: ``I am going to sequence these pictures\ldots'' invariably means they are going to be arranged into a type of order, not necessarily a sequence.

Much of the confusion between series and sequence, two explicitly separate structures, comes from a muddling of terms. Sequence is defined in the dictionary as series, and vice versa. Even those who work with ordering pictures often use the two terms interchangeably:

\begin{quote}
``There is no more surprising, yet, in its naturalness and organic sequence, simpler form than the photographic series. This is the logical culmination of photography/vision in motion. The series is no longer a `picture' and the canons of pictorial aesthetics can only be applied to it \textit{mutatis mutandis}. Here the single picture loses its separate identity and becomes a structural element of the related whole which is the thing in itself. In this sequence of separate but inseparable parts, a photographic series---photographic comics, pamphlets, books---can be either a potent weapon or tender poetry. But first must come to the realization that the knowledge of photography is just as important as that of the alphabet. The illiterate of the future will be the person ignorant of the use of the camera as well as of the pen.''

---Maholy Nagy, published in \textit{Telehor}, 1936
\end{quote}

So, there are impressive precedents for using the two terms freely. To do so makes it impossible to differentiate between two different and specific types of organization---a series and a sequence. Semantically, I would not care if series and sequence were used interchangeably, as long as the two distinct phenomena are reckoned with. Perhaps they could be called a simple series/sequence versus a complex series/sequence, or a 2-D series/sequence versus a 3-D series/sequence.

I have pondered over calling a series a contiguous referral, and a sequence a contingental referral; it is easier simply to say series, and sequence. But in using either of these two terms, one must be very careful that it is preceded with a definition.

I have found it prudent whenever possible not to use the terms group, series, and sequence. Instead, I refer to them by what they do:

\begin{itemize}
  \item A group is a static collection.
  \item A series is a constructed linking movement.
  \item Based on cause and effect, a sequence is a constructed conditional movement.
\end{itemize}

\rule{\textwidth}{0.4pt}

\subsection*{Glossary}

\textbf{Cadence}: ``The natural rhythm of language determined by its inherent alternation of stressed and unstressed syllables. When more precisely used in verse, the term cadence refers to the arrangement of the rhythms of speech into highly organized patterns.''

\textbf{Contrapuntal}: see counterpoint.

\textbf{Counterpoint}: two or more melodic lines sounding simultaneously, as opposed to harmony. ``A melody added to a melody as accompaniment. The art of plural melody.'' A book employs counterpoint by interweaving.

\textbf{Group}: 1) One of three types of units. The only connective in a group is a common denominator of resemblance or relationship. Motif is stated, but can never be evolved except by context: each additional one clarifies the common denominator. 2) A listing, not modified. There is no movement except leaf flow.

\textbf{Harmony}: chords. The simultaneous occurrence of musical tones as opposed to melody, a succession of tones. See texture, counterpoint.

\textbf{Melody}: a succession of musical tones as opposed to harmony (tones sounded simultaneously: chords). See texture.

\textbf{Leaf}: the front and back of a page. Recto/verso. A sheet.

\textbf{Leaf Flow}: 1) unordered mechanical movement from the front to the back cover of a book, through the pages with no detours. 2) pagination.

\textbf{Meter}: ``rhythmical structure concerned with the division into measures consisting of a uniform number of beats or time units.'' See prosody.

\textbf{Pacing}: 1) The modulation of time in a book. 2) Cadence.

\textbf{Progression}: movement through a book which is directed as opposed to leaf flow.
\begin{itemize}
  \item In a \textbf{series}, arithmetical progression is in a straight line, proceeding at a constant difference, as opposed to geometric (constant factor). The constancy, subject matter, and narration is modified by referral and pacing.
  \item In a \textbf{sequence}, geometric progression is causal, in and out of a picture, and from picture to picture.
\end{itemize}

\textbf{Prosody}: 1) ``The systematic study of metrical structure, including varieties of poetic feet and meters, rhymes and rhyming patterns, types of stanzas and strophes, and fixed forms.'' 2) ``Verse can be based on accent or quantity. Meter refers to the pattern of stressed and unstressed syllables. Meter is allowed to vary while retaining its basic pattern, to avoid becoming mechanical. Or, the regularity of meter can be abandoned for cadence, which uses rhythm, approximating the flow of speech.''
\begin{itemize}
  \item Meters in English verse: iambic, trochaic, anapestic, dactylic, spondaic, pyrrhic
  \item Metrical units, or feet, are in a line and numbered: 1 monometer, 2 dimeter, 3 trimeter, 4 tetrameter, 5 pentameter, 6 hexameter, 7 heptameter, 8 octameter
\end{itemize}

\textbf{Random Referral}: undirected movement by free association made by the viewer in any unit. See referral.

\textbf{Recollection}: 1) Relating by referral to a previous element in a series. 2) An aspect of serial movement, along with linkage-forward and preview.

\textbf{Recto/Verso}: the two sides of a leaf in a codex. Recto is the right-hand, verso is the back of that same leaf, not the page facing the recto.

\textbf{(Direct) Referral}: a relationship made between pictures by the bookmaker.
\begin{itemize}
  \item In a \textbf{series}, it is limited to adjacent pictures causing linkage-forward, supplemented by occasional non-adjacent references of recollection and preview.
  \item In a \textbf{sequence}, the reference is causal, from one element (subject matter, form, composition, etc.) to any other element within that or any other picture. Movement is multidirectional, within or between pictures.
  \item Reference can extend beyond a unit by symbolism, innuendo, double entendre, narration, and parable. See random referral.
\end{itemize}

\textbf{Revelation}: 1) The act of turning the page to experience the book. 2) One of the five elements of the book.

\textbf{Rhythm}
\begin{enumerate}
  \item The counterpart of motion.
  \item ``The flow of cadences in written or spoken language. The regular rise and fall of sounds, whether in pitch, stress, speed\ldots Attention to quantities of syllables, accents and pauses\ldots Flow of movement which groups by recurrent heavy and light accent, \ldots grouping beats into measures, as 4/4 rhythm. Movement marked by regular recurrence of regular alternation in features, elements, phenomena, etc. Hence, periodicity.''
\end{enumerate}

\textbf{Sequence}: one of three types of units. It differs from a series by the elaborate references within or between any other picture. Instead of a strong linear progression of narrative, movement is causal.

\textbf{Sequential Movement}: 1) A geometric progression by cause and effect, as opposed to linkage of serial movement. 2) A planar texture.

\textbf{Serial Movement}
\begin{enumerate}
  \item Linkage-forward, with possible references of preview and recollection.
  \item A linear progression, the means of narration and metamorphosis.
\end{enumerate}

\textbf{Series}: One of three types of units; the others are group and sequence.

\textbf{Texture}
\begin{enumerate}
  \item The combination of the vertical harmony with the horizontal movement of melody.
  \item A visual interweaving: In a \textbf{series}, by means of plot and sub-plot. In a \textbf{sequence}, by the layered (vertical) references of elements with the horizontal movement of direct referral: sequential movement.
  \item One of the formal elements.
\end{enumerate}

\textbf{Thematic Development}
\begin{itemize}
  \item In a \textbf{series}, statement and the occurring variations or ensuing metamorphosis. The unveiling of plot. Narration.
  \item In a \textbf{sequence}, the use of motif as structure. The evolution and modulation of an idea through sequential movement.
\end{itemize}

\textbf{Theme}: Motif which is chosen as a major idea to be evolved.

\end{document}
